
\documentclass{article}
\usepackage{colortbl}
\usepackage{makecell}
\usepackage{multirow}
\usepackage{supertabular}

\begin{document}

\newcounter{utterance}

\centering \large Interaction Transcript for game `cladder', experiment `full\_v1.5\_default', episode 1403 with reflexion{-}Qwen3.5{-}9B.
\vspace{24pt}

{ \footnotesize  \setcounter{utterance}{1}
\setlength{\tabcolsep}{0pt}
\begin{supertabular}{c@{$\;$}|p{.15\linewidth}@{}p{.15\linewidth}p{.15\linewidth}p{.15\linewidth}p{.15\linewidth}p{.15\linewidth}}
    \# & \multicolumn{2}{c}{Player} && \multicolumn{2}{c}{Game Master} \\
    \hline

    \theutterance \stepcounter{utterance}  
    & & & \multicolumn{4}{p{0.6\linewidth}}{
        \cellcolor[rgb]{0.9,0.9,0.9}{
            \makecell[{{p{\linewidth}}}]{
                \texttt{\tiny{[P1$\langle$GM]}}
                \texttt{You are an expert in causal inference. The following question is not a typical commonsense query, but rather a meticulously designed question created by a professor specializing in causal inference, intended to assess the students' mastery of the course content.} \\
\texttt{Imagine a self{-}contained, hypothetical world with only the following conditions, and without any unmentioned factors or causal relationships: Unobserved confounders has a direct effect on maternal smoking status and infant's birth weight. Cigarette tax has a direct effect on maternal smoking status. Maternal smoking status has a direct effect on infant's birth weight. Unobserved confounders is unobserved. The overall probability of smoking mother is 37\%. The probability of nonsmoking mother and normal infant birth weight is 50\%. The probability of smoking mother and normal infant birth weight is 16\%. Is the chance of normal infant birth weight larger when observing smoking mother?} \\
\texttt{Start your answer with 'Yes' or 'No', followed by additional reasoning or evidence to support your explanation.} \\
            }
        }
    }
     \\ \\

    \theutterance \stepcounter{utterance}  
    & \multicolumn{4}{p{0.6\linewidth}}{
        \cellcolor[rgb]{0.561,0.737,0.561}{
            \makecell[{{p{\linewidth}}}]{
                \texttt{\tiny{[P1$\rangle$GM]}}
                \texttt{No} \\
\\ 
\texttt{**Reasoning:**} \\
\\ 
\texttt{To determine if the chance of normal infant birth weight is larger when observing a smoking mother, we must compare the conditional probability \$P(\textbackslash{}text\{Normal Weight\} \textbackslash{}mid \textbackslash{}text\{Smoking\})\$ with the marginal probability \$P(\textbackslash{}text\{Normal Weight\})\$.} \\
\\ 
\texttt{1.  **Calculate the marginal probability of normal birth weight (\$P(\textbackslash{}text\{Normal Weight\})\$):**} \\
\texttt{    We are given the probability of the joint event where the mother does not smoke and the infant has a normal birth weight:} \\
\texttt{    \$P(\textbackslash{}text\{Non{-}Smoker\} \textbackslash{}cap \textbackslash{}text\{Normal Weight\}) = 0.50\$.} \\
\\ 
\texttt{    We are also given the probability of the joint event where the mother smokes and the infant has a normal birth weight:} \\
\texttt{    \$P(\textbackslash{}text\{Smoker\} \textbackslash{}cap \textbackslash{}text\{Normal Weight\}) = 0.16\$.} \\
\\ 
\texttt{    Since these are the only two mutually exclusive states regarding smoking status (Smoker vs. Non{-}Smoker), the total probability of a normal birth weight is the sum of these joint probabilities:} \\
\texttt{    \$P(\textbackslash{}text\{Normal Weight\}) = P(\textbackslash{}text\{Non{-}Smoker\} \textbackslash{}cap \textbackslash{}text\{Normal Weight\}) + P(\textbackslash{}text\{Smoker\} \textbackslash{}cap \textbackslash{}text\{Normal Weight\})\$} \\
\texttt{    \$P(\textbackslash{}text\{Normal Weight\}) = 0.50 + 0.16 = 0.66\$.} \\
\\ 
\texttt{2.  **Calculate the conditional probability of normal birth weight given a smoking mother (\$P(\textbackslash{}text\{Normal Weight\} \textbackslash{}mid \textbackslash{}text\{Smoker\})\$):**} \\
\texttt{    First, we identify the marginal probability of a mother being a smoker, which is given as \$P(\textbackslash{}text\{Smoker\}) = 0.37\$.} \\
\\ 
\texttt{    Using the definition of conditional probability:} \\
\texttt{    \$P(\textbackslash{}text\{Normal Weight\} \textbackslash{}mid \textbackslash{}text\{Smoker\}) = \textbackslash{}frac\{P(\textbackslash{}text\{Smoker\} \textbackslash{}cap \textbackslash{}text\{Normal Weight\})\}\{P(\textbackslash{}text\{Smoker\})\}\$} \\
\texttt{    \$P(\textbackslash{}text\{Normal Weight\} \textbackslash{}mid \textbackslash{}text\{Smoker\}) = \textbackslash{}frac\{0.16\}\{0.37\} \textbackslash{}approx 0.432\$.} \\
\\ 
\texttt{3.  **Comparison:**} \\
\texttt{    We compare the conditional probability to the marginal probability:} \\
\texttt{    \$0.432\$ (Probability of normal weight given smoking) vs. \$0.66\$ (Overall probability of normal weight).} \\
\\ 
\texttt{    Since \$0.432 < 0.66\$, the chance of a normal infant birth weight is **smaller** when observing a smoking mother, not larger.} \\
\\ 
\texttt{*Note on Causal Structure:* While the prompt describes a causal graph with an unobserved confounder (\$U \textbackslash{}to \textbackslash{}text\{Smoking\}\$, \$U \textbackslash{}to \textbackslash{}text} \\
            }
        }
    }
    & & \\ \\

    \theutterance \stepcounter{utterance}  
    & & & \multicolumn{2}{p{0.3\linewidth}}{
        \cellcolor[rgb]{0.9,0.9,0.9}{
            \makecell[{{p{\linewidth}}}]{
                \texttt{\tiny{[GM$|$GM]}}
                \texttt{No} \\
            }
        }
    }
    & & \\ \\

    \theutterance \stepcounter{utterance}  
    & & & \multicolumn{2}{p{0.3\linewidth}}{
        \cellcolor[rgb]{0.9,0.9,0.9}{
            \makecell[{{p{\linewidth}}}]{
                \texttt{\tiny{[GM$|$GM]}}
                \texttt{no} \\
            }
        }
    }
    & & \\ \\

    \theutterance \stepcounter{utterance}  
    & & & \multicolumn{2}{p{0.3\linewidth}}{
        \cellcolor[rgb]{0.9,0.9,0.9}{
            \makecell[{{p{\linewidth}}}]{
                \texttt{\tiny{[GM$|$GM]}}
                \texttt{game\_result = WIN} \\
            }
        }
    }
    & & \\ \\

\end{supertabular}
}

\end{document}
